\documentclass[11pt]{article}
\usepackage{amsmath,amssymb,amsthm}
\usepackage[margin=1in]{geometry}
\newtheorem{proposition}{Proposition}
\begin{document}
\title{Cell C4: complete proofs for five and six lines}
\author{}
\date{26 September 2026}
\maketitle

\noindent\fbox{\parbox{\dimexpr\linewidth-2\fboxsep-2\fboxrule}{\textbf{Official statement (Cell 4).} ``Two concrete instances of the next case, $N = d+2$: five lines in three dimensions and six lines in four. In each, the conjectured optimum repeats two of the coordinate axes. Prove that any 5 lines in $\mathbb R^3$ satisfy $S \le 4\pi$, and that any 6 lines in $\mathbb R^4$ satisfy $S \le 13\pi/2$.''}}
\medskip

For unit representatives $x_1,\ldots,x_N\in\mathbb R^d$, put
$g_{ij}=\langle x_i,x_j\rangle$,
$A=\sum_{i<j}\arcsin|g_{ij}|$, and
$S=\sum_{i<j}\arccos|g_{ij}|$.  Since
$\theta_{ij}=\pi/2-\arcsin|g_{ij}|$, the two requested inequalities
are both equivalent to $A\geq\pi$.

\section*{Elementary lemmas}

\begin{proposition}[Planar triples and quadruples]
For three lines in a plane, the sum of their pairwise acute angles
is at most $\pi$.  For four lines in a plane, it is at most $2\pi$.
Equivalently, their respective $A$-values are at least $\pi/2$
and $\pi$.
\end{proposition}
\begin{proof}
Represent the three lines by three points on a circle of circumference
$\pi$; acute angle is shortest circular distance.  If the three
successive gaps are all at most $\pi/2$, the three pair distances sum
to the circumference $\pi$.  If one gap is greater than $\pi/2$,
the sum of pair distances is twice the sum of the other two gaps,
which is less than $\pi$.  This proves the triple statement.
For four lines, sum the triple statement over their four triples.
Each pair appears twice, so twice the four-line angle sum is at most
$4\pi$.
For later use, averaging the three-line $A\geq\pi/2$ bound over
the ten triples among five planar lines shows that their $A$-value
is at least $5\pi/3$: every pair occurs in three triples.
\end{proof}

\begin{proposition}[Orthogonal basis]
If $e_1,\ldots,e_d$ are an orthonormal basis and $z$ is a unit
vector, then $\sum_{j=1}^d\arcsin|\langle z,e_j\rangle|\geq\pi/2$.
\end{proposition}
\begin{proof}
For $a,b\geq0$ with $a^2+b^2\leq1$ one has
\[
 \arcsin a+\arcsin b\geq\arcsin\sqrt{a^2+b^2}.
\]
If the left side is already at least $\pi/2$, this is immediate.
Otherwise, taking sines reduces the claim to
$\sqrt{(1-a^2)(1-b^2)}\geq ab$, which follows from
$a^2+b^2\leq1$.  Combine the $d$ coordinates successively and use
$\sum_j|\langle z,e_j\rangle|^2=1$.
\end{proof}

\begin{proposition}[Orthogonality saturation]
Fix $N$ lines in $\mathbb R^d$ and let $S$ be their total acute-angle
sum.  There exists a maximizer of $S$ such that, for every line
$\ell_i$, either it is orthogonal to all other lines or its
orthogonality neighbors span a space of dimension at least $d-1$.
\end{proposition}
\begin{proof}
The product of projective spheres is compact, so a maximizer exists.
Among all maximizers choose one with the largest possible number of
orthogonal pairs.  Let $W$ be the span of the representatives of the
lines orthogonal to $\ell_i$.  Suppose $\dim W\leq d-2$ and that
$\ell_i$ is not orthogonal to all the others.  Then $W^\perp$ has
dimension at least two, and there is a great circle in $W^\perp$
through a unit representative $x_i$ on which the inner product with
some non-neighbor is nonconstant.  Vary only $x_i$ on this circle.

On an interval with no new orthogonality event, each term
$\arccos|\langle x_i(t),x_j\rangle|$ is convex.  To see this,
write its inner product, up to sign, as $r\cos(t-\phi)$, with
$0\leq r\leq1$ and $\cos(t-\phi)>0$.  For $r<1$ its second
derivative is
\[
 \frac{r(1-r^2)\cos(t-\phi)}
 {(1-r^2\cos^2(t-\phi))^{3/2}}\geq0.
\]
For $r=1$ the function is piecewise linear with a convex corner
at coincidence.  Terms with $r=0$ are constant.  Thus the sum is
convex between consecutive orthogonality events.  One endpoint of
the event interval containing $x_i$ has value at least its initial
value.  Since the initial configuration was a global maximizer, the
endpoint is another maximizer.  All old orthogonality pairs involving
$i$ persist, since the circle lies in $W^\perp$, and at least one
new pair appears.  This contradicts the choice of maximizer.
\end{proof}

The same argument has a strict form that we will use only to describe
equality.  At \emph{any} maximizer, suppose an orthogonality-neighbor
span $W$ has dimension at most $d-2$.  On every two-dimensional circle
$U\subset W^\perp$ through $x_i$, a non-neighbor $y$ with
$0<\|\operatorname{proj}_U y\|<1$ would make its angle term strictly
convex between events, contradicting maximality.  Thus every
non-neighbor belongs to every such $U$.  If $\dim W^\perp=2$, all
non-neighbors lie in $W^\perp$; if $\dim W^\perp\geq3$, they lie on
$\ell_i$, the intersection of these circle planes.  Hence the
configuration splits into orthogonal subfamilies, unless $W=0$
and all lines coincide.  The latter is not a maximizer for $N\geq2$
and $d\geq2$.

\section*{Five lines in three dimensions}

\begin{proposition}
For any five lines in $\mathbb R^3$, $A\geq\pi$, or equivalently
$S\leq4\pi$.
\end{proposition}
\begin{proof}
It suffices to prove the bound for a maximizing configuration chosen
as in the orthogonality-saturation proposition.  Let $H$ be its graph
of orthogonal pairs.  If a vertex is adjacent to all four others,
those four lines lie in a plane; their angle sum is at most $2\pi$,
and the four angles involving the universal vertex total $2\pi$.

Otherwise the neighbors of every vertex span a plane, so $H$ has
minimum degree at least two.  If $H$ contains a triangle, its three
lines form an orthonormal basis.  The orthogonal-basis proposition
applied to each of the remaining two lines gives $A\geq\pi$.

If $H$ is triangle-free, minimum degree two on five vertices forces
$H=C_5$ or $H=K_{2,3}$.  Indeed, degree two everywhere gives $C_5$;
if some vertex has degree three, its three neighbors are mutually
nonadjacent, and the fifth vertex must meet all three to give them
degree at least two, giving $K_{2,3}$.  In the latter case the two
parts span orthogonal subspaces.  If the part of size two consists
of independent lines, the three lines in the other part coincide
and contribute $3\pi/2$ to $A$.  Otherwise the two lines coincide,
contributing $\pi/2$, while the other three are planar and
contribute at least another $\pi/2$.

It remains to treat $C_5$.  Label its edges
$12,23,34,45,51$.  Choose coordinates with $x_1=e_1$ and
$x_2=e_2$.  Since $x_2\perp x_3$ and $x_1\perp x_5$, write,
up to signs irrelevant to the absolute inner products,
\[
 x_3=(\cos t,0,\sin t),\qquad
 x_5=(0,\cos s,\sin s),\qquad 0\leq t,s\leq\pi/2.
\]
The line $x_4$ is perpendicular to both $x_3,x_5$.  Put
$D=\sqrt{1-\sin^2t\sin^2s}$.  If $D=0$, then $x_3=x_5=e_3$
as lines and $x_1,x_2,x_3$ are an orthonormal basis, a case already
handled.  Otherwise, the five nonzero-potential pairs are the
nonedges of $C_5$, and direct cross-product calculation gives
\[
 A=\pi-t-s+u(t,s)+v(t,s)+\gamma(t,s)
  =\pi+F(t,s),
 \quad F(t,s):=u(t,s)+v(t,s)+\gamma(t,s)-t-s,
\]
where
\[
 u=\arctan(\tan t\cos s),\qquad
 v=\arctan(\tan s\cos t),\qquad
 \gamma=\arcsin(\sin t\sin s).
\]
The expression is interpreted continuously at endpoints.  We show
$F\geq0$.  For fixed $s\in[0,\pi/2)$, differentiation gives
\[
 \frac{\partial F}{\partial t}
 =\frac{\cos s}{1+\sin t\sin s}
  +\frac{\cos t\sin s}{\sqrt{1-\sin^2t\sin^2s}}-1.
\]
Both nonconstant terms decrease with $t$.  For the second,
\[
 \frac{d}{dt}\left(\frac{\cos t}{D}\right)
 =-\frac{\sin t\cos^2s}{D^3}\leq0.
\]
Hence $F$ is concave in $t$.
But $F(0,s)=F(\pi/2,s)=0$, so $F(t,s)\geq0$.
The case $s=\pi/2$ follows by continuity.  Therefore $A\geq\pi$.
\end{proof}

\section*{Six lines in four dimensions}

\begin{proposition}
Any four lines in $\mathbb R^3$ satisfy $A\geq\pi/2$.
\end{proposition}
\begin{proof}
Choose a maximizing configuration for their acute-angle sum with
the largest number of orthogonal pairs.  If one line is orthogonal
to the other three, the latter are planar, so their $A$-value is
at least $\pi/2$.  Otherwise saturation gives minimum degree two
in the orthogonality graph on four vertices.  If it contains a
triangle, that triangle is an orthonormal basis and the fourth line
contributes at least $\pi/2$ by the orthogonal-basis proposition.
If it is triangle-free, it must be $K_{2,2}$: the two parts span
orthogonal spaces, and since their dimensions sum to at most three,
at least one part is a pair of coincident lines, contributing
$\pi/2$ to $A$.
\end{proof}

\begin{proposition}
For any six lines in $\mathbb R^4$, $A\geq\pi$, or equivalently
$S\leq13\pi/2$.
\end{proposition}
\begin{proof}
The same saturation proposition gives a useful restriction.  At a
maximizing configuration of six lines in $\mathbb R^4$, a universal
orthogonality vertex reduces the bound to the five-line theorem above:
the other five lines lie in its orthogonal $\mathbb R^3$, while its
five incident angles total $5\pi/2$.  Four mutually orthogonal
lines also settle the case, by applying the orthogonal-basis
proposition to each of the other two lines.

If neither occurs, let $J$ be the graph of nonorthogonal pairs.
Every orthogonality neighborhood spans a three-dimensional space,
so every vertex of $J$ has degree at most two.  If $J$ is connected,
it is a path or a cycle on six vertices.  A path is impossible:
the Gram matrix, in path order, is an irreducible symmetric
tridiagonal matrix, whose eigenvalues are simple (the first
coordinate of an eigenvector determines all the others by the
three-term recurrence).  Its nullity is therefore at most one,
whereas six vectors in $\mathbb R^4$ have Gram nullity at least two.
If $J$ is disconnected, its components are mutually orthogonal
families.  For any component of at most five lines with span rank
$r$, the preceding planar bounds, the four-line proposition, and
the five-line theorem show directly that its $A$-value is at least
$(n-r)\pi/2$: check $r=1$ by counting coincident pairs, $r=2$
by the planar bounds and the five-line triple average, and
$r=3$ by the four- and five-line theorems.  A component of rank
$4$ leaves no room for any other nonempty orthogonal component.
Summing over components gives $A\geq(6-4)\pi/2=\pi$.
It remains to handle $J=C_6$.
This case also follows from the general cycle lemma in \texttt{p1\_c5.tex}.
The two alternating triples of its
vertices are each pairwise orthogonal.  Denote their unit vectors by
$u_1,u_2,u_3$ and $v_1,v_2,v_3$, relabeling so that
$u_i\perp v_i$ for $i=1,2,3$ and the six other cross pairs are the
edges of $J$.  Complete $u_1,u_2,u_3$ to an orthonormal basis of
$\mathbb R^4$, so $u_i=e_i$.  Orienting coordinate axes and lines
does not change $A$.  We may therefore write
\[
 v_1=(0,C,Pc,Ps),\qquad
 v_2=(D,0,Ws,-Wc),
\]
where $P=\sin p$, $C=\cos p$, $W=\sin w$, $D=\cos w$,
$c=\cos q$, $s=\sin q$, and $p,w,q\in[0,\pi/2]$.
Indeed, the first expression parameterizes a unit vector orthogonal
to $e_1$; the second parameterizes a unit vector orthogonal to
both $e_2$ and $v_1$.  For the strict $C_6$ pattern all six cross
inner products are nonzero, so $C,D,c,s>0$.  Boundary cases follow
by continuity or have already been covered by a graph with more
orthogonality edges.

The vector $v_3$ is orthogonal to $e_3,v_1,v_2$.  Solving these
three linear equations and normalizing shows that the absolute
values of its first two coordinates are
\[
 X=\frac{WCc}{H},\qquad Y=\frac{PDs}{H},\qquad
 H^2=C^2D^2+P^2D^2s^2+W^2C^2c^2.
\]
The six nonzero-potential pair terms are consequently
\[
 A=\arcsin C+\arcsin D+\arcsin(Pc)+\arcsin(Ws)
       +\arcsin X+\arcsin Y
   =\pi+F(q),
\]
where
\[
 F(q)=\arcsin(Pc)+\arcsin(Ws)+\arcsin X+\arcsin Y-p-w.
\]
We prove $F(q)\geq0$.  Its endpoint values are
$F(0)=F(\pi/2)=0$.  Set
\[
 a=\sqrt{1-P^2c^2},\qquad b=\sqrt{1-W^2s^2},\qquad
 M=\frac{CD}{H^2}.
\]
Direct differentiation, using $1-X^2=D^2a^2/H^2$ and
$1-Y^2=C^2b^2/H^2$, gives
\[
 F'(q)=\frac{c(W+PM)}{b}-\frac{s(P+WM)}{a}.
\]
For $0<p,w<\pi/2$, the sign of $F'(q)$ is the sign of $R(q)-1$,
where
\[
 R(q)=\frac{c}{s}\frac{a}{b}\frac{W+PM}{P+WM}.
\]
The first two factors have a nonincreasing product: multiplying
its logarithmic derivative by $sc$ yields
\[
 -1+\frac{P^2s^2c^2}{a^2}
    +\frac{W^2s^2c^2}{b^2}\leq0,
\]
because $a^2=C^2+P^2s^2\geq P^2s^2$ and
$b^2=D^2+W^2c^2\geq W^2c^2$.  The last factor is also
nonincreasing.  Indeed $H^2=D^2+(W^2-P^2)c^2$, so
\[
 M'=\frac{2CD(W^2-P^2)sc}{H^4},
\]
and the logarithmic derivative of $(W+PM)/(P+WM)$ is
\[
 \frac{M'(P^2-W^2)}{(W+PM)(P+WM)}
 =-\frac{2CD(W^2-P^2)^2sc}
 {H^4(W+PM)(P+WM)}\leq0.
\]
Thus $R$ decreases from $+\infty$ at $q=0$ to $0$ at
$q=\pi/2$.  The derivative $F'$ changes sign only from positive
to negative, so $F$ has its minimum at an endpoint.  Both endpoint
values are zero; hence $A\geq\pi$.  Cases with $p$ or $w$ on the
boundary follow by continuity.  This completes the six-line proof.
\end{proof}

\noindent The six-line case also follows from the cycle lemma of Cell~5 (applied with $m=6$),
which gives a second, independent proof.

\section*{Sharpness}

Three planar lines with directions $0,\pi/3,2\pi/3$ have three
pairwise acute angles $\pi/3$, hence contribute $\pi/2$ to $A$.
Four planar lines with directions $0,\pi/2,\delta,\delta+\pi/2$
have $A=\pi$ for every $0\leq\delta\leq\pi/2$.
For five lines in $\mathbb R^3$, put the planar triple in one
two-dimensional summand and repeat the perpendicular line twice;
their contributions are $\pi/2+\pi/2=\pi$.  Alternatively, use
the planar quadruple and one perpendicular line.  For six lines in
$\mathbb R^4$, put one such planar triple in each of two mutually
orthogonal planes.  Their contributions again sum to $\pi$.
Thus both upper bounds are attained.

\section*{Equality configurations}

Here is an exact description, with repeated lines allowed.  Let $D$
be a pair of coincident lines; let $T$ be any three lines in a plane
whose three successive gaps on the projective circle are all at most
$\pi/2$; let $Q$ be any four lines in a plane that are the union of
two orthogonal pairs.  The symbol $+$ below denotes an orthogonal
direct sum of the spans of the indicated blocks, and $1$ denotes a
single line.  Up to an orthogonal transformation and relabeling,
all equality configurations are
\[
 \begin{array}{c|c}
 (N,d)&\text{blocks}\\ \hline
 (5,3)&Q+1\quad\text{or}\quad T+D\\
 (6,4)&Q+1+1\quad\text{or}\quad T+D+1
                  \quad\text{or}\quad T+T.
 \end{array}
\]
Each displayed family attains $A=\pi$, by the planar bounds with
equality and by $A(D)=\pi/2$.

For completeness, we justify that the list is exhaustive.  Equality
in the orthogonal-basis proposition holds precisely when the extra
vector has at most two nonzero basis coordinates: combining two
positive coordinates is strict whenever their squared sum is less
than one.  Equality for a planar triple is exactly the gap condition
in the definition of $T$.  For four cyclically ordered planar lines,
write the four consecutive gaps as $a,b,c,d$, with
$a+b+c+d=\pi$.  Equality in the four-line bound forces equality in
each of its four triple bounds, hence
\[
 a+b,\ b+c,\ c+d,\ d+a\leq\pi/2.
\]
The four left sides sum to $2\pi$, so all are $\pi/2$; opposite
lines form two orthogonal pairs.  This is exactly $Q$.

An equality configuration is a global maximizer.  If it is not an
orthogonal sum, the strict form of orthogonality saturation gives
the same minimum-degree conclusions for its own orthogonality graph
as those used above.  For four lines in $\mathbb R^3$, its graph is
then a triangle or $K_{2,2}$; the latter is an orthogonal sum, and
equality in the former puts the fourth line in a coordinate plane,
leaving a basis line orthogonal to the other three.  Thus every
four-line equality case is $T+1$.

For five lines in $\mathbb R^3$, the $K_{2,3}$ case is an orthogonal
sum and the exact $C_5$ case is strict: with $0<t,s<\pi/2$, the
function $F(t,s)$ above is strictly concave in $t$ and is positive
between its zero endpoints.  In the orthogonal-basis case, the two
extra lines each have support of size at most two and are orthogonal
to each other.  If their supports are disjoint or equal, they split
off an orthogonal coordinate block; supports meeting in exactly one
coordinate have nonzero inner product.  Hence
this case also splits orthogonally.  Decomposing repeatedly, and
using the sharp component bounds in ranks one, two, and three,
leaves exactly $Q+1$ or $T+D$.

For six lines in $\mathbb R^4$, the exact $C_6$ case is strict because
$R(q)$ is strictly decreasing in the interior, so $F(q)>0$ there.
The path is impossible and a disconnected nonorthogonality graph
already gives an orthogonal sum.  If four lines form an orthonormal
basis, equality requires the other two to be orthogonal and each
to have support in at most two basis coordinates.  Their supports
are either disjoint or equal; supports meeting in only one coordinate
would have nonzero inner product.  Thus this case
also splits.  Finally, the universal-vertex case splits by definition.
Every equality case is therefore an orthogonal sum of components
of sizes at most five.  The component inequalities used above imply
that equality requires total span rank four and equality in each
component.  The only positive-defect block types needed are
$D$ (defect one), $T$ (defect one), and $Q$ (defect two); the
four-line and five-line equality classifications just proved split
all other blocks further.  A zero-defect component at equality has
$A=0$, so its lines split into orthogonal singletons.  Counting
six lines and rank four gives
exactly the three displayed families.


\section*{An independent short estimate}

\begin{proposition}
For $(N,d)=(5,3)$ or $(6,4)$, the requested inequality holds if
$|g_{ij}|\leq 1/2$ for every $i\ne j$.  In the first case the stronger
bound $A\geq 10\pi/9$ holds.
\end{proposition}
\begin{proof}
The function $f(t)=\arcsin(t)/t^2$ decreases on $(0,1/2]$.  Indeed,
the sign of $f'(t)$ is the sign of
$t/\sqrt{1-t^2}-2\arcsin t$, which is negative there because
$t/\sqrt{1-t^2}\leq 2t/\sqrt3<2t\leq2\arcsin t$.
Consequently,
\[
  \arcsin t\geq \frac{2\pi}{3}t^2
  \qquad (0\leq t\leq 1/2).
\]
Let $T=\sum_i x_i x_i^{\mathsf T}$.  It is positive semidefinite,
$\operatorname{tr}T=N$, and Cauchy--Schwarz on its at most $d$
nonzero eigenvalues gives
\[
 N+2\sum_{i<j}g_{ij}^2
 =\operatorname{tr}(T^2)\geq \frac{N^2}{d}.
\]
Thus
\[
 A\geq\frac{2\pi}{3}\sum_{i<j}g_{ij}^2
 \geq \frac{\pi N(N-d)}{3d}.
\]
This last expression is $10\pi/9$ for $(N,d)=(5,3)$ and
$\pi$ for $(N,d)=(6,4)$.
\end{proof}

The hypothesis is exactly that every acute line angle is at least
$\pi/3$.  This gives a separate check on that part of the
configuration space.

\end{document}
